% FAC-EN-U0051
% French source: fac_body.tex lines 1830--1857.
% Authority: PDF physical pp. 50--51; printed pp. 245--246 (no. 53 complete).

\medskip
\phantomsection\label{fac:III:1:no:53}
\textbf{53. Cohomology of irreducible algebraic curves.}
If $V$ is an irreducible algebraic variety of dimension $1$, the closed
subsets of $V$ distinct from $V$ are precisely the \textit{finite} subsets.
If $F$ is a finite subset of $V$ and $x$ is a point of $F$, put
\[
V_x^F=(V\setminus F)\cup\{x\}.
\]
The sets $V_x^F$, $x\in F$, form a finite open covering
\phantomsection\label{fac:III:1:no:53:cover:B-F}
$\mathfrak{B}^F$ of $V$.

\phantomsection\label{fac:III:1:lemma:3}
\textsc{Lemma 3.}
\textit{The coverings $\mathfrak{B}^F$ of the preceding type are
arbitrarily fine.}

Let $\mathfrak{U}=\{U_i\}_{i\in I}$ be an open covering of $V$. We may
suppose it finite, since $V$ is quasi-compact, and we may also suppose that
$U_i\neq\varnothing$ for every $i\in I$. Put $F_i=V\setminus U_i$. Each
$F_i$ is finite, and so is
\[
F=\bigcup_{i\in I}F_i.
\]
We show that $\mathfrak{B}^F<\mathfrak{U}$, which will prove the lemma. Let
$x\in F$. Since the $U_i$ cover $V$, there is an $i\in I$ such that
$x\notin F_i$. As $F\supset F_i$, we then have
$F\setminus\{x\}\supset F_i$; equivalently, $V_x^F\subset U_i$. Thus
$\mathfrak{B}^F<\mathfrak{U}$, as required.

\phantomsection\label{fac:III:1:lemma:4}
\textsc{Lemma 4.}
\textit{Let $\mathcal{F}$ be a sheaf on $V$, and let $F$ be a finite subset
of $V$. Then
\[
H^n(\mathfrak{B}^F,\mathcal{F})=0
\]
for $n\geq2$.}

Put $W=V\setminus F$. If $x_0,\ldots,x_n$ are distinct and $n\geq1$, then
\[
V_{x_0}^F\cap\cdots\cap V_{x_n}^F=W.
\]
Consequently, if $G=\Gamma(W,\mathcal{F})$, the alternating complex
$C'(\mathfrak{B}^F,\mathcal{F})$ is isomorphic, in dimensions $\geq1$, to
$C'(S(F),G)$, where $S(F)$ denotes the simplex whose vertex set is $F$.
It follows that
\phantomsection\label{fac:III:1:no:53:formula:simplex-cohomology}
\[
H^n(\mathfrak{B}^F,\mathcal{F})=H^n(S(F),G)=0
\qquad\text{for }n\geq2,
\]
since the cohomology of a simplex is trivial.

\hyperref[fac:III:1:lemma:3]{Lemmas~3} \hyperref[fac:III:1:lemma:4]{and~4} plainly imply:

\phantomsection\label{fac:III:1:prop:4}
\textsc{Proposition 4.}
\textit{If $V$ is an irreducible algebraic curve and $\mathcal{F}$ is any
sheaf on $V$, then
\[
H^n(V,\mathcal{F})=0\qquad\text{for }n\geq2.
\]}

\phantomsection\label{fac:III:1:no:53:remark}
\textsc{Remark.}
I do not know whether a result analogous to the \hyperref[fac:III:1:prop:4]{preceding one} is valid for
varieties of arbitrary dimension.
